% СГЕНЕРИРОВАНО инструменты/generate.py — руками не править

\pfDefGlava{Ч1}{Ch1}{Дроби, проценты, единицы}{Числа и вычисления}
\pfGlavaTemy{Ч1}{Ч1.1,Ч1.2,Ч1.3,Ч1.4,Ч1.5}
\pfDefGlava{А1}{A1}{Преобразования выражений}{Алгебра}
\pfGlavaTemy{А1}{А1.1,А1.2,А1.3}
\pfDefGlava{А2}{A2}{Степени, корни, логарифмы}{Алгебра}
\pfGlavaTemy{А2}{А2.1,А2.2,А2.3,А2.4,А2.5,А2.6}
\pfDefGlava{А3}{A3}{Уравнения}{Алгебра}
\pfGlavaTemy{А3}{А3.1,А3.2,А3.3,А3.4,А3.5,А3.6,А3.7,А3.8,А3.9}
\pfDefGlava{А4}{A4}{Прогрессии}{Алгебра}
\pfGlavaTemy{А4}{А4.1,А4.2}
\pfDefGlava{Т1}{T1}{Основы тригонометрии}{Тригонометрия}
\pfGlavaTemy{Т1}{Т1.1,Т1.2,Т1.3,Т1.4,Т1.5,Т1.6}
\pfDefGlava{Ф1}{F1}{Функция и её свойства}{Функции и графики}
\pfGlavaTemy{Ф1}{Ф1.1,Ф1.2,Ф1.3,Ф1.4,Ф1.5}
\pfDefGlava{Ф2}{F2}{Элементарные функции}{Функции и графики}
\pfGlavaTemy{Ф2}{Ф2.1,Ф2.2,Ф2.3,Ф2.4,Ф2.5,Ф2.6,Ф2.7,Ф2.8}
\pfDefGlava{А5}{A5}{Неравенства}{Алгебра}
\pfGlavaTemy{А5}{А5.1,А5.2,А5.3,А5.4,А5.5,А5.6,А5.7,А5.8}
\pfDefGlava{Т2}{T2}{Тригонометрические формулы}{Тригонометрия}
\pfGlavaTemy{Т2}{Т2.1,Т2.2,Т2.3,Т2.4,Т2.5,Т2.6,Т2.7,Т2.8}
\pfDefGlava{Т3}{T3}{Тригонометрические уравнения}{Тригонометрия}
\pfGlavaTemy{Т3}{Т3.1,Т3.2,Т3.3,Т3.4,Т3.5,Т3.6,Т3.7}
\pfDefGlava{Д1}{D1}{Производная}{Производная}
\pfGlavaTemy{Д1}{Д1.1,Д1.2,Д1.3,Д1.4,Д1.5}
\pfDefGlava{Д2}{D2}{Исследование функции}{Производная}
\pfGlavaTemy{Д2}{Д2.1,Д2.2,Д2.3,Д2.4,Д2.5}
\pfDefGlava{П1}{P1}{Углы и прямые}{Планиметрия}
\pfGlavaTemy{П1}{П1.1,П1.2}
\pfDefGlava{П2}{P2}{Треугольник}{Планиметрия}
\pfGlavaTemy{П2}{П2.1,П2.2,П2.3,П2.4,П2.5,П2.6,П2.7,П2.8,П2.9}
\pfDefGlava{П3}{P3}{Прямоугольный треугольник}{Планиметрия}
\pfGlavaTemy{П3}{П3.1,П3.2,П3.3,П3.4,П3.5,П3.6}
\pfDefGlava{П4}{P4}{Четырёхугольники}{Планиметрия}
\pfGlavaTemy{П4}{П4.1,П4.2,П4.3,П4.4}
\pfDefGlava{П5}{P5}{Окружность}{Планиметрия}
\pfGlavaTemy{П5}{П5.1,П5.2,П5.3,П5.4,П5.5,П5.6,П5.7,П5.8,П5.9,П5.10,П5.11}
\pfDefGlava{П6}{P6}{Вписанные и описанные окружности}{Планиметрия}
\pfGlavaTemy{П6}{П6.1,П6.2,П6.3,П6.4,П6.5,П6.6,П6.7,П6.8}
\pfDefGlava{П7}{P7}{Метрические соотношения в треугольнике}{Планиметрия}
\pfGlavaTemy{П7}{П7.1,П7.2,П7.3,П7.4,П7.5,П7.6}
\pfDefGlava{П8}{P8}{Площади}{Планиметрия}
\pfGlavaTemy{П8}{П8.1,П8.2,П8.3,П8.4,П8.5}
\pfDefGlava{П9}{P9}{Ортоцентр и высоты}{Планиметрия}
\pfGlavaTemy{П9}{П9.1,П9.2}
\pfDefGlava{К1}{K1}{Векторы на плоскости}{Векторы и координаты}
\pfGlavaTemy{К1}{К1.1,К1.2,К1.8,К1.3,К1.4,К1.5,К1.6,К1.7}
\pfDefGlava{С1}{S1}{Многогранники и тела вращения}{Стереометрия}
\pfGlavaTemy{С1}{С1.1,С1.2,С1.3,С1.4,С1.5,С1.6,С1.7,С1.8,С1.9,С1.10}
\pfDefGlava{С2}{S2}{Взаимное расположение в пространстве}{Стереометрия}
\pfGlavaTemy{С2}{С2.1,С2.2,С2.3,С2.4,С2.5,С2.6}
\pfDefGlava{С3}{S3}{Углы и расстояния}{Стереометрия}
\pfGlavaTemy{С3}{С3.1,С3.2,С3.3,С3.4,С3.5,С3.6,С3.8,С3.7}
\pfDefGlava{С4}{S4}{Сечения}{Стереометрия}
\pfGlavaTemy{С4}{С4.1}
\pfDefGlava{К2}{K2}{Координатный метод в пространстве}{Векторы и координаты}
\pfGlavaTemy{К2}{К2.1,К2.2,К2.3,К2.4,К2.5,К2.6,К2.7,К2.8,К2.9}
\pfDefGlava{В1}{V1}{Комбинаторика}{Вероятность и статистика}
\pfGlavaTemy{В1}{В1.2,В1.1}
\pfDefGlava{В2}{V2}{Вероятность}{Вероятность и статистика}
\pfGlavaTemy{В2}{В2.1,В2.2,В2.3,В2.4,В2.5,В2.6,В2.7,В2.8,В2.9}
\pfDefGlava{В3}{V3}{Случайные величины}{Вероятность и статистика}
\pfGlavaTemy{В3}{В3.1,В3.2,В3.3,В3.4}
\pfDefGlava{М1}{M1}{Текстовые и прикладные задачи}{Модели: текстовые и экономические задачи}
\pfGlavaTemy{М1}{М1.1,М1.2,М1.7,М1.3,М1.4,М1.5,М1.6}
\pfDefGlava{М2}{M2}{Финансовая математика и оптимизация}{Модели: текстовые и экономические задачи}
\pfGlavaTemy{М2}{М2.1,М2.2,М2.3,М2.4,М2.5,М2.6}
\pfDefGlava{Ф3}{F3}{Графики на плоскости (для параметра)}{Функции и графики}
\pfGlavaTemy{Ф3}{Ф3.1,Ф3.2,Ф3.3,Ф3.4,Ф3.5,Ф3.6}
\pfDefGlava{Ч2}{Ch2}{Делимость и целые числа}{Числа и вычисления}
\pfGlavaTemy{Ч2}{Ч2.1,Ч2.2,Ч2.3,Ч2.4,Ч2.5,Ч2.6,Ч2.7,Ч2.8,Ч2.9,Ч2.10,Ч2.11,Ч2.12}

\pfDefTema{Ч1.1}{Ch1-1}{Дроби: основное свойство, перевод десятичных дробей в обыкновенные}{Ч1}{}{20}{стр. 216 (старое задание 19)}{}
\pfDefTema{Ч1.2}{Ch1-2}{Среднее арифметическое}{Ч1}{}{20}{стр. 216 (старое задание 19)}{}
\pfDefTema{Ч1.3}{Ch1-3}{Проценты: процент от числа, число по проценту, процентное отношение}{Ч1}{\tref{Ч1.1}}{11}{стр. 91–92 (старое задание 10)}{\rref{Р11-01}}
\pfDefTema{Ч1.4}{Ch1-4}{Увеличение и уменьшение на p\% (множитель 1 ± p/100)}{Ч1}{\tref{Ч1.3}}{13}{стр. 152 (старое задание 16)}{}
\pfDefTema{Ч1.5}{Ch1-5}{Перевод единиц измерения (скорость, время, длина, площадь, объём)}{Ч1}{}{11}{стр. 80; 92 (старое задание 9; 10)}{}
\pfDefTema{Ч2.1}{Ch2-1}{Чётность и её свойства}{Ч2}{}{20}{стр. 208 (старое задание 19)}{}
\pfDefTema{Ч2.2}{Ch2-2}{Признаки делимости}{Ч2}{}{20}{стр. 208 (старое задание 19)}{}
\pfDefTema{Ч2.3}{Ch2-3}{Десятичная запись числа}{Ч2}{}{20}{стр. 209 (старое задание 19)}{}
\pfDefTema{Ч2.4}{Ch2-4}{Последняя цифра числа}{Ч2}{\tref{Ч2.3}}{20}{стр. 210 (старое задание 19)}{}
\pfDefTema{Ч2.5}{Ch2-5}{Простые и составные числа, основная теорема арифметики}{Ч2}{}{20}{стр. 210 (старое задание 19)}{}
\pfDefTema{Ч2.6}{Ch2-6}{Количество делителей числа}{Ч2}{\tref{Ч2.5}}{20}{стр. 209 (старое задание 19)}{}
\pfDefTema{Ч2.7}{Ch2-7}{НОД и НОК, взаимно простые числа}{Ч2}{\tref{Ч2.5}}{20}{стр. 210 (старое задание 19)}{}
\pfDefTema{Ч2.8}{Ch2-8}{Свойства НОД и НОК}{Ч2}{\tref{Ч2.7}}{20}{стр. 211 (старое задание 19)}{}
\pfDefTema{Ч2.9}{Ch2-9}{Деление с остатком, алгебра остатков}{Ч2}{}{20}{стр. 212 (старое задание 19)}{}
\pfDefTema{Ч2.10}{Ch2-10}{Сравнения по модулю}{Ч2}{\tref{Ч2.9}}{20}{стр. 212 (старое задание 19)}{}
\pfDefTema{Ч2.11}{Ch2-11}{Уравнения в целых числах}{Ч2}{\tref{Ч2.1}, \tref{Ч2.5}}{20}{стр. 211 (старое задание 19)}{}
\pfDefTema{Ч2.12}{Ch2-12}{Метод «оценка + пример»}{Ч2}{\tref{Ч2.11}}{20}{— нет в методичке}{}
\pfDefTema{А1.1}{A1-1}{Формулы сокращённого умножения; вынесение общего множителя и группировка}{А1}{}{7, 8, 14, 16}{стр. 49 (старое задание 6)}{}
\pfDefTema{А1.2}{A1-2}{Модуль: определение, геометрический смысл, свойства, базовые переходы}{А1}{}{8, 16, 19}{стр. 146–150 (старое задание 15)}{\rref{Р16-01}}
\pfDefTema{А1.3}{A1-3}{Многочлены: теорема Безу, подбор корня, деление уголком, схема Горнера}{А1}{\tref{А1.1}}{14, 16}{стр. 150–151 (старое задание 15)}{}
\pfDefTema{А2.1}{A2-1}{Степень с целым и рациональным показателем, свойства степеней}{А2}{}{7, 8}{стр. 58; 65 (старое задание 6; 7)}{}
\pfDefTema{А2.2}{A2-2}{Корни (арифметический квадратный, кубический) и их свойства}{А2}{\tref{А2.1}}{8}{стр. 65 (старое задание 7)}{}
\pfDefTema{А2.3}{A2-3}{Логарифм: определение и ограничения}{А2}{\tref{А2.1}}{7}{стр. 61 (старое задание 6)}{}
\pfDefTema{А2.4}{A2-4}{Свойства логарифмов}{А2}{\tref{А2.3}}{7, 8}{стр. 62; 66 (старое задание 6; 7)}{}
\pfDefTema{А2.5}{A2-5}{Логарифмы чётных степеней и произведений: модули в формулах}{А2}{\tref{А2.4}, \tref{А1.2}}{16}{стр. 145 (старое задание 15)}{}
\pfDefTema{А2.6}{A2-6}{Сравнение чисел с корнями, логарифмами, дробями}{А2}{\tref{А2.2}, \tref{А2.4}}{14}{стр. 117–119 (старое задание 13)}{}
\pfDefTema{А3.1}{A3-1}{Линейные и квадратные уравнения, дискриминант}{А3}{\tref{А1.1}}{7}{стр. 49–50 (старое задание 6)}{}
\pfDefTema{А3.2}{A3-2}{Неполные квадратные уравнения}{А3}{\tref{А3.1}}{7}{стр. 50–51 (старое задание 6)}{}
\pfDefTema{А3.3}{A3-3}{Теорема Виета (прямая и обратная)}{А3}{\tref{А3.1}}{7}{стр. 51–52 (старое задание 6)}{}
\pfDefTema{А3.4}{A3-4}{Уравнения вида (f(x))³ = a}{А3}{\tref{А2.2}}{7}{стр. 53 (старое задание 6)}{}
\pfDefTema{А3.5}{A3-5}{Рациональные уравнения}{А3}{\tref{А3.1}}{7}{стр. 53–54 (старое задание 6)}{}
\pfDefTema{А3.6}{A3-6}{Иррациональные уравнения}{А3}{\tref{А2.2}, \tref{А3.1}}{7, 19}{стр. 54–58; 201 (старое задание 6; 18)}{}
\pfDefTema{А3.7}{A3-7}{Показательные уравнения}{А3}{\tref{А2.1}}{7}{стр. 58–60; 64 (старое задание 6)}{}
\pfDefTema{А3.8}{A3-8}{Логарифмические уравнения}{А3}{\tref{А2.4}}{7, 19}{стр. 61–64; 201 (старое задание 6; 18)}{}
\pfDefTema{А3.9}{A3-9}{Уравнения с модулем}{А3}{\tref{А1.2}}{19}{стр. 201 (старое задание 18)}{}
\pfDefTema{А4.1}{A4-1}{Арифметическая прогрессия}{А4}{}{11, 13, 20}{— нет в методичке}{}
\pfDefTema{А4.2}{A4-2}{Геометрическая прогрессия}{А4}{\tref{А2.1}}{11}{— нет в методичке}{}
\pfDefTema{А5.1}{A5-1}{Метод интервалов}{А5}{\tref{А3.1}}{16}{стр. 140 (старое задание 15)}{}
\pfDefTema{А5.2}{A5-2}{Квадратные неравенства}{А5}{\tref{А5.1}}{10}{стр. 80 (старое задание 9)}{}
\pfDefTema{А5.3}{A5-3}{ОДЗ, ограничения и равносильные переходы (сводная тема)}{А5}{\tref{А2.2}, \tref{А2.3}}{16}{стр. 141 (старое задание 15)}{}
\pfDefTema{А5.4}{A5-4}{Иррациональные неравенства: равносильные переходы}{А5}{\tref{А5.1}, \tref{А3.6}}{16}{стр. 142 (старое задание 15)}{}
\pfDefTema{А5.5}{A5-5}{Показательные неравенства, метод рационализации}{А5}{\tref{А2.1}, \tref{Ф2.6}}{16}{стр. 142–144 (старое задание 15)}{}
\pfDefTema{А5.6}{A5-6}{Логарифмические неравенства, метод рационализации}{А5}{\tref{А2.4}, \tref{Ф2.7}}{16}{стр. 143–144 (старое задание 15)}{}
\pfDefTema{А5.7}{A5-7}{Неравенства с модулем}{А5}{\tref{А1.2}, \tref{А5.1}}{16}{стр. 146–150 (старое задание 15)}{\rref{Р16-01}}
\pfDefTema{А5.8}{A5-8}{Системы и совокупности неравенств}{А5}{\tref{А5.1}}{16}{стр. 151 (старое задание 15)}{}
\pfDefTema{Т1.1}{T1-1}{sin, cos, tg, ctg острого угла в прямоугольном треугольнике}{Т1}{}{1}{стр. 6 (старое задание 1)}{}
\pfDefTema{Т1.2}{T1-2}{Радианная мера угла}{Т1}{}{7, 8}{— нет в методичке}{}
\pfDefTema{Т1.3}{T1-3}{Тригонометрическая окружность, знаки по четвертям}{Т1}{\tref{Т1.2}}{8}{стр. 70 (старое задание 7)}{}
\pfDefTema{Т1.4}{T1-4}{Таблица значений тригонометрических функций}{Т1}{\tref{Т1.3}}{1, 8}{стр. 7; 67 (старое задание 1; 7)}{}
\pfDefTema{Т1.5}{T1-5}{Чётность тригонометрических функций}{Т1}{\tref{Т1.3}}{8}{стр. 69 (старое задание 7)}{}
\pfDefTema{Т1.6}{T1-6}{Периодичность тригонометрических функций}{Т1}{\tref{Т1.3}}{8, 14}{— нет в методичке}{}
\pfDefTema{Т2.1}{T2-1}{Основное тригонометрическое тождество и следствия}{Т2}{\tref{Т1.3}}{8}{стр. 68 (старое задание 7)}{}
\pfDefTema{Т2.2}{T2-2}{Формулы сложения}{Т2}{\tref{Т2.1}}{8}{стр. 68 (старое задание 7)}{}
\pfDefTema{Т2.3}{T2-3}{Формулы двойного угла}{Т2}{\tref{Т2.2}}{8}{стр. 69 (старое задание 7)}{}
\pfDefTema{Т2.4}{T2-4}{Формулы понижения степени}{Т2}{\tref{Т2.3}}{8}{стр. 69 (старое задание 7)}{}
\pfDefTema{Т2.5}{T2-5}{Преобразование суммы и разности в произведение}{Т2}{\tref{Т2.2}}{8}{стр. 68 (старое задание 7)}{}
\pfDefTema{Т2.6}{T2-6}{Универсальная тригонометрическая подстановка}{Т2}{\tref{Т2.3}}{8}{стр. 68 (старое задание 7)}{}
\pfDefTema{Т2.7}{T2-7}{Формулы приведения}{Т2}{\tref{Т1.3}}{8}{стр. 70–71 (старое задание 7)}{}
\pfDefTema{Т2.8}{T2-8}{Метод вспомогательного аргумента}{Т2}{\tref{Т2.2}}{14, 19}{стр. 115; 198 (старое задание 13; 18)}{}
\pfDefTema{Т3.1}{T3-1}{Обратные тригонометрические функции: arcsin, arccos, arctg, arcctg}{Т3}{\tref{Т1.3}}{7, 14}{— нет в методичке}{}
\pfDefTema{Т3.2}{T3-2}{Уравнение sin x = a}{Т3}{\tref{Т3.1}}{14}{стр. 106–107 (старое задание 13)}{}
\pfDefTema{Т3.3}{T3-3}{Уравнение cos x = a}{Т3}{\tref{Т3.1}}{14}{стр. 107 (старое задание 13)}{}
\pfDefTema{Т3.4}{T3-4}{Уравнения tg x = a и ctg x = a}{Т3}{\tref{Т3.1}}{14}{стр. 108–109 (старое задание 13)}{}
\pfDefTema{Т3.5}{T3-5}{Отбор корней на промежутке (окружность, неравенство, таблица)}{Т3}{\tref{Т3.2}, \tref{Т3.3}}{14}{стр. 110–112 (старое задание 13)}{}
\pfDefTema{Т3.6}{T3-6}{Методы решения: замена, разложение на множители, группировка}{Т3}{\tref{Т3.2}, \tref{Т2.3}}{14}{стр. 113–114 (старое задание 13)}{}
\pfDefTema{Т3.7}{T3-7}{Однородные уравнения I и II степени}{Т3}{\tref{Т3.4}}{14}{стр. 114–115 (старое задание 13)}{}
\pfDefTema{Ф1.1}{F1-1}{Координатная плоскость, оси}{Ф1}{}{12}{стр. 93 (старое задание 11)}{}
\pfDefTema{Ф1.2}{F1-2}{Монотонность функции; единственность корня}{Ф1}{\tref{Ф1.1}}{19}{стр. 196–197 (старое задание 18)}{}
\pfDefTema{Ф1.3}{F1-3}{Чётность и симметрия функции}{Ф1}{}{19}{стр. 205–207 (старое задание 18)}{}
\pfDefTema{Ф1.4}{F1-4}{Область значений функции, оценка}{Ф1}{\tref{Ф1.2}}{19}{стр. 198–200 (старое задание 18)}{}
\pfDefTema{Ф1.5}{F1-5}{Преобразования графиков: сдвиги, отражения, модуль}{Ф1}{\tref{Ф1.1}}{19}{стр. 201–203 (старое задание 18)}{}
\pfDefTema{Ф2.1}{F2-1}{Линейная функция; k = tg α}{Ф2}{\tref{Ф1.1}, \tref{Т1.1}}{12}{стр. 93–94 (старое задание 11)}{}
\pfDefTema{Ф2.2}{F2-2}{Квадратичная функция, вершина параболы}{Ф2}{\tref{А3.1}}{12}{стр. 94–95 (старое задание 11)}{}
\pfDefTema{Ф2.3}{F2-3}{Расположение корней квадратного трёхчлена}{Ф2}{\tref{Ф2.2}}{19}{стр. 190–191 (старое задание 18)}{}
\pfDefTema{Ф2.4}{F2-4}{Гипербола, асимптоты}{Ф2}{\tref{Ф1.5}}{12}{стр. 96 (старое задание 11)}{}
\pfDefTema{Ф2.5}{F2-5}{Функция корня}{Ф2}{\tref{А2.2}}{12}{стр. 97 (старое задание 11)}{}
\pfDefTema{Ф2.6}{F2-6}{Показательная функция}{Ф2}{\tref{А2.1}}{12}{стр. 98 (старое задание 11)}{}
\pfDefTema{Ф2.7}{F2-7}{Логарифмическая функция}{Ф2}{\tref{А2.3}}{12}{стр. 99 (старое задание 11)}{}
\pfDefTema{Ф2.8}{F2-8}{Взаимное расположение графиков, точки пересечения}{Ф2}{\tref{Ф2.1}, \tref{Ф2.2}}{12}{стр. 100–101 (старое задание 11)}{}
\pfDefTema{Ф3.1}{F3-1}{Уравнения фигур на плоскости: окружность, «ромб», отрезок}{Ф3}{\tref{К1.2}}{19}{стр. 191 (старое задание 18)}{}
\pfDefTema{Ф3.2}{F3-2}{Построение множеств с модулями (раскрытие по областям)}{Ф3}{\tref{А1.2}, \tref{Ф1.5}}{19}{стр. 191–195 (старое задание 18)}{}
\pfDefTema{Ф3.3}{F3-3}{Метод областей}{Ф3}{\tref{А5.1}}{19}{стр. 203–204 (старое задание 18)}{}
\pfDefTema{Ф3.4}{F3-4}{Траектории движения особых точек графика}{Ф3}{\tref{Ф1.5}}{19}{стр. 202–203 (старое задание 18)}{}
\pfDefTema{Ф3.5}{F3-5}{Графический метод в плоскости (x; a)}{Ф3}{\tref{Ф3.2}}{19}{— нет в методичке}{}
\pfDefTema{Ф3.6}{F3-6}{Условия касания: прямая и окружность, прямая и парабола}{Ф3}{\tref{Ф3.1}, \tref{Д1.5}}{19}{— нет в методичке}{}
\pfDefTema{Д1.1}{D1-1}{Производная, геометрический смысл}{Д1}{\tref{Ф2.1}}{9}{стр. 72–73 (старое задание 8)}{}
\pfDefTema{Д1.2}{D1-2}{Физический смысл производной}{Д1}{\tref{Д1.1}}{9}{— нет в методичке}{}
\pfDefTema{Д1.3}{D1-3}{Таблица производных}{Д1}{\tref{Д1.1}}{17}{стр. 102–103 (старое задание 12)}{}
\pfDefTema{Д1.4}{D1-4}{Правила дифференцирования}{Д1}{\tref{Д1.3}}{17}{стр. 103 (старое задание 12)}{}
\pfDefTema{Д1.5}{D1-5}{Касательная к графику функции}{Д1}{\tref{Д1.1}}{9}{стр. 74–75 (старое задание 8)}{}
\pfDefTema{Д2.1}{D2-1}{Знак производной и монотонность}{Д2}{\tref{Д1.1}, \tref{Ф1.2}}{9, 17}{стр. 75–76; 104 (старое задание 8; 12)}{}
\pfDefTema{Д2.2}{D2-2}{Экстремумы: необходимое и достаточное условия}{Д2}{\tref{Д2.1}}{9, 17}{стр. 77–78; 104 (старое задание 8; 12)}{}
\pfDefTema{Д2.3}{D2-3}{Чтение графиков функции и производной}{Д2}{\tref{Д2.2}}{9}{стр. 76–79 (старое задание 8)}{}
\pfDefTema{Д2.4}{D2-4}{Алгоритм поиска точек экстремума}{Д2}{\tref{Д2.2}, \tref{А5.1}}{17}{стр. 104 (старое задание 12)}{}
\pfDefTema{Д2.5}{D2-5}{Наибольшее и наименьшее значение на отрезке}{Д2}{\tref{Д2.4}}{17}{стр. 105 (старое задание 12)}{}
\pfDefTema{П1.1}{P1-1}{Смежные и вертикальные углы}{П1}{}{1}{стр. 5 (старое задание 1)}{}
\pfDefTema{П1.2}{P1-2}{Параллельные прямые: углы при секущей}{П1}{\tref{П1.1}}{1, 18}{— нет в методичке}{}
\pfDefTema{П2.1}{P2-1}{Сумма углов треугольника, внешний угол}{П2}{\tref{П1.2}}{1}{стр. 4 (старое задание 1)}{}
\pfDefTema{П2.2}{P2-2}{Биссектриса, медиана, высота}{П2}{}{1}{стр. 8 (старое задание 1)}{}
\pfDefTema{П2.3}{P2-3}{Признаки равенства треугольников}{П2}{}{1}{стр. 12 (старое задание 1)}{}
\pfDefTema{П2.4}{P2-4}{Равнобедренный треугольник}{П2}{\tref{П2.3}}{1, 18}{— нет в методичке}{}
\pfDefTema{П2.5}{P2-5}{Теорема Фалеса, обобщённая, обратные}{П2}{\tref{П1.2}}{18}{стр. 170–172 (старое задание 17)}{}
\pfDefTema{П2.6}{P2-6}{Средняя линия треугольника}{П2}{\tref{П2.5}}{1}{стр. 9 (старое задание 1)}{}
\pfDefTema{П2.7}{P2-7}{Подобие треугольников: признаки, отношение площадей}{П2}{\tref{П2.5}}{1}{стр. 13–14 (старое задание 1)}{}
\pfDefTema{П2.8}{P2-8}{Свойство биссектрисы внутреннего и внешнего угла}{П2}{\tref{П2.7}}{18}{стр. 166 (старое задание 17)}{}
\pfDefTema{П2.9}{P2-9}{Точка пересечения медиан (деление 2 : 1)}{П2}{\tref{П2.6}}{1, 18}{— нет в методичке}{}
\pfDefTema{П3.1}{P3-1}{Теорема Пифагора и обратная}{П3}{\tref{П2.1}}{1, 18}{стр. 16; 163 (старое задание 1; 17)}{}
\pfDefTema{П3.2}{P3-2}{Катет против угла 30°}{П3}{\tref{П3.1}}{18}{стр. 162 (старое задание 17)}{}
\pfDefTema{П3.3}{P3-3}{Медиана, проведённая к гипотенузе}{П3}{\tref{П2.2}}{18}{стр. 163 (старое задание 17)}{}
\pfDefTema{П3.4}{P3-4}{Высота к гипотенузе: подобие и пропорциональные отрезки}{П3}{\tref{П2.7}, \tref{П3.1}}{15, 18}{стр. 126; 163 (старое задание 14; 17)}{}
\pfDefTema{П3.5}{P3-5}{Соответственные элементы подобных прямоугольных треугольников}{П3}{\tref{П3.4}}{18}{стр. 164 (старое задание 17)}{}
\pfDefTema{П3.6}{P3-6}{Высота, медиана и биссектриса из прямого угла}{П3}{\tref{П3.3}, \tref{П3.4}}{18}{стр. 164–165 (старое задание 17)}{}
\pfDefTema{П4.1}{P4-1}{Параллелограмм, ромб, прямоугольник, квадрат: определения и свойства}{П4}{\tref{П2.3}, \tref{П1.2}}{18}{стр. 162 (старое задание 17)}{}
\pfDefTema{П4.2}{P4-2}{Трапеция: определение и свойства}{П4}{\tref{П4.1}}{1, 18}{— нет в методичке}{}
\pfDefTema{П4.3}{P4-3}{Средняя линия трапеции; отрезок между серединами диагоналей}{П4}{\tref{П2.6}, \tref{П4.2}}{1}{стр. 10–11 (старое задание 1)}{}
\pfDefTema{П4.4}{P4-4}{Сумма углов многоугольника, правильные многоугольники (сноска)}{П4}{\tref{П4.1}}{1}{— нет в методичке}{}
\pfDefTema{П5.1}{P5-1}{Окружность и круг: длина, площадь}{П5}{}{1}{стр. 21 (старое задание 1)}{}
\pfDefTema{П5.2}{P5-2}{Центральный и вписанный углы}{П5}{\tref{П5.1}}{1}{стр. 25–27 (старое задание 1)}{}
\pfDefTema{П5.3}{P5-3}{Вписанный угол, опирающийся на диаметр; отрезок, видимый под углом}{П5}{\tref{П5.2}}{18}{стр. 179 (старое задание 17)}{}
\pfDefTema{П5.4}{P5-4}{Касательная: перпендикулярность радиусу, отрезки касательных}{П5}{\tref{П5.1}}{1, 18}{— нет в методичке}{}
\pfDefTema{П5.5}{P5-5}{Угол между касательной и хордой}{П5}{\tref{П5.2}, \tref{П5.4}}{18}{стр. 177 (старое задание 17)}{}
\pfDefTema{П5.6}{P5-6}{Углы между хордами, секущими и касательными}{П5}{\tref{П5.2}}{18}{стр. 178 (старое задание 17)}{}
\pfDefTema{П5.7}{P5-7}{Хорды и дуги}{П5}{\tref{П5.2}}{18}{стр. 179–181 (старое задание 17)}{}
\pfDefTema{П5.8}{P5-8}{Пересекающиеся хорды; касательная и секущая}{П5}{\tref{П2.7}, \tref{П5.4}}{18}{стр. 181–182 (старое задание 17)}{}
\pfDefTema{П5.9}{P5-9}{Площадь сектора и сегмента}{П5}{\tref{П5.1}, \tref{Т1.2}}{18}{стр. 177 (старое задание 17)}{}
\pfDefTema{П5.10}{P5-10}{Взаимное расположение двух окружностей}{П5}{\tref{П5.4}}{18}{стр. 186–187 (старое задание 17)}{}
\pfDefTema{П5.11}{P5-11}{Касающиеся окружности: общая касательная, подобие}{П5}{\tref{П5.10}}{18}{стр. 187–188 (старое задание 17)}{}
\pfDefTema{П6.1}{P6-1}{Вписанная окружность треугольника}{П6}{\tref{П2.2}, \tref{П5.4}}{1}{стр. 22–23 (старое задание 1)}{}
\pfDefTema{П6.2}{P6-2}{Четырёхугольник, описанный около окружности}{П6}{\tref{П6.1}}{1}{стр. 23 (старое задание 1)}{}
\pfDefTema{П6.3}{P6-3}{Описанная окружность треугольника}{П6}{\tref{П5.3}}{1, 18}{стр. 24; 163 (старое задание 1; 17)}{}
\pfDefTema{П6.4}{P6-4}{Вписанный четырёхугольник: признаки, теорема Птолемея}{П6}{\tref{П5.2}}{1, 18}{стр. 24; 176 (старое задание 1; 17)}{}
\pfDefTema{П6.5}{P6-5}{Радиус вписанной окружности прямоугольного треугольника}{П6}{\tref{П6.1}, \tref{П3.1}}{1, 18}{стр. 22; 164 (старое задание 1; 17)}{}
\pfDefTema{П6.6}{P6-6}{Равносторонний треугольник: элементы, радиусы, площадь}{П6}{\tref{П6.1}, \tref{П6.3}}{1}{стр. 17–18 (старое задание 1)}{}
\pfDefTema{П6.7}{P6-7}{Вневписанная окружность}{П6}{\tref{П6.1}}{18}{стр. 165; 178 (старое задание 17)}{}
\pfDefTema{П6.8}{P6-8}{Лемма о трезубце}{П6}{\tref{П6.7}, \tref{П5.2}}{18}{стр. 188–189 (старое задание 17)}{}
\pfDefTema{П7.1}{P7-1}{Теорема косинусов}{П7}{\tref{П3.1}, \tref{Т1.3}, \tref{Т2.7}}{1}{стр. 15 (старое задание 1)}{}
\pfDefTema{П7.2}{P7-2}{Теорема синусов}{П7}{\tref{П6.3}, \tref{П5.2}}{1}{стр. 15 (старое задание 1)}{}
\pfDefTema{П7.3}{P7-3}{Длина медианы}{П7}{\tref{П7.1}}{18}{стр. 166 (старое задание 17)}{}
\pfDefTema{П7.4}{P7-4}{Длина биссектрисы}{П7}{\tref{П2.8}}{18}{стр. 167 (старое задание 17)}{}
\pfDefTema{П7.5}{P7-5}{Углы между биссектрисами; деление биссектрис точкой пересечения}{П7}{\tref{П2.8}}{18}{стр. 167–168 (старое задание 17)}{}
\pfDefTema{П7.6}{P7-6}{Теоремы Менелая и Чевы}{П7}{\tref{П2.7}}{18}{стр. 169–170 (старое задание 17)}{}
\pfDefTema{П8.1}{P8-1}{Площадь треугольника: все формулы}{П8}{\tref{П7.2}, \tref{П6.1}, \tref{Т2.7}}{1}{стр. 19 (старое задание 1)}{}
\pfDefTema{П8.2}{P8-2}{Площади четырёхугольников}{П8}{\tref{П4.1}, \tref{П4.2}}{1}{стр. 20 (старое задание 1)}{}
\pfDefTema{П8.3}{P8-3}{Отношения площадей треугольников}{П8}{\tref{П8.1}}{18}{стр. 172–174 (старое задание 17)}{}
\pfDefTema{П8.4}{P8-4}{Равновеликие треугольники («рельсы Евклида»)}{П8}{\tref{П8.1}}{18}{стр. 162 (старое задание 17)}{}
\pfDefTema{П8.5}{P8-5}{Площадь через радиусы вневписанных окружностей}{П8}{\tref{П6.7}, \tref{П8.1}}{18}{стр. 165–166 (старое задание 17)}{}
\pfDefTema{П9.1}{P9-1}{Высоты и вписанные четырёхугольники; ортотреугольник}{П9}{\tref{П6.4}}{18}{стр. 174–175 (старое задание 17)}{}
\pfDefTema{П9.2}{P9-2}{Ортоцентр и описанная окружность}{П9}{\tref{П9.1}, \tref{П6.3}}{18}{стр. 182–186 (старое задание 17)}{}
\pfDefTema{К1.1}{K1-1}{Вектор, координаты вектора}{К1}{\tref{Ф1.1}}{2}{стр. 28–29 (старое задание 2)}{}
\pfDefTema{К1.2}{K1-2}{Длина вектора, расстояние между точками}{К1}{\tref{К1.1}, \tref{П3.1}}{2}{стр. 30 (старое задание 2)}{}
\pfDefTema{К1.3}{K1-3}{Сложение и вычитание векторов}{К1}{\tref{К1.1}}{2}{стр. 30–31 (старое задание 2)}{}
\pfDefTema{К1.4}{K1-4}{Умножение вектора на число}{К1}{\tref{К1.1}}{2}{стр. 32 (старое задание 2)}{}
\pfDefTema{К1.5}{K1-5}{Коллинеарные векторы}{К1}{\tref{К1.4}}{2}{стр. 33 (старое задание 2)}{}
\pfDefTema{К1.6}{K1-6}{Скалярное произведение, угол между векторами, перпендикулярность векторов}{К1}{\tref{К1.2}, \tref{П7.1}}{2}{стр. 34 (старое задание 2)}{}
\pfDefTema{К1.7}{K1-7}{Расстояние от точки до прямой на плоскости}{К1}{\tref{К1.6}}{19}{стр. 195 (старое задание 18)}{}
\pfDefTema{К1.8}{K1-8}{Координаты середины отрезка}{К1}{\tref{К1.2}}{2}{— нет в методичке}{}
\pfDefTema{К2.1}{K2-1}{Координаты точки и вектора в пространстве}{К2}{\tref{К1.1}, \tref{С1.1}}{15}{— нет в методичке}{}
\pfDefTema{К2.2}{K2-2}{Деление отрезка в данном отношении}{К2}{\tref{К2.1}}{15}{стр. 139 (старое задание 14)}{}
\pfDefTema{К2.3}{K2-3}{Уравнение плоскости, вектор нормали}{К2}{\tref{К2.1}, \tref{К1.6}}{15}{стр. 133–134 (старое задание 14)}{}
\pfDefTema{К2.4}{K2-4}{Угол между прямыми через координаты}{К2}{\tref{К2.1}, \tref{К1.6}}{15}{стр. 133 (старое задание 14)}{}
\pfDefTema{К2.5}{K2-5}{Угол между прямой и плоскостью через координаты}{К2}{\tref{К2.3}, \tref{С3.2}}{15}{стр. 135 (старое задание 14)}{}
\pfDefTema{К2.6}{K2-6}{Угол между плоскостями через координаты}{К2}{\tref{К2.3}, \tref{С3.3}}{15}{стр. 135 (старое задание 14)}{}
\pfDefTema{К2.7}{K2-7}{Расстояние от точки до плоскости}{К2}{\tref{К2.3}}{15}{стр. 135 (старое задание 14)}{}
\pfDefTema{К2.8}{K2-8}{Расстояние между скрещивающимися прямыми и между параллельными плоскостями}{К2}{\tref{К2.7}, \tref{С3.7}}{15}{стр. 136–137 (старое задание 14)}{}
\pfDefTema{К2.9}{K2-9}{Доказательства методом координат}{К2}{\tref{К2.3}}{15}{стр. 137–139 (старое задание 14)}{}
\pfDefTema{С1.1}{S1-1}{Куб и прямоугольный параллелепипед}{С1}{\tref{П3.1}}{3}{стр. 35 (старое задание 3)}{}
\pfDefTema{С1.2}{S1-2}{Прямая призма}{С1}{\tref{П8.1}}{3}{стр. 35 (старое задание 3)}{}
\pfDefTema{С1.3}{S1-3}{Пирамида; правильная пирамида: проекция вершины в центр основания, апофема}{С1}{\tref{П8.1}}{3}{стр. 35 (старое задание 3)}{}
\pfDefTema{С1.4}{S1-4}{Цилиндр}{С1}{\tref{П5.1}}{3}{стр. 35 (старое задание 3)}{}
\pfDefTema{С1.5}{S1-5}{Конус}{С1}{\tref{П5.1}, \tref{П3.1}}{3}{стр. 36 (старое задание 3)}{}
\pfDefTema{С1.6}{S1-6}{Шар}{С1}{}{3}{стр. 36 (старое задание 3)}{}
\pfDefTema{С1.7}{S1-7}{Комбинации тел}{С1}{\tref{С1.1}, \tref{С1.4}, \tref{С1.5}, \tref{С1.6}}{3}{стр. 36–37 (старое задание 3)}{}
\pfDefTema{С1.8}{S1-8}{Подобие тел (k, k², k³)}{С1}{\tref{П2.7}}{3}{стр. 38–39 (старое задание 3)}{}
\pfDefTema{С1.9}{S1-9}{Отношение объёмов частей тела (отсечённая пирамида, часть конуса)}{С1}{\tref{С1.3}, \tref{С1.5}, \tref{С1.8}}{3}{— нет в методичке}{}
\pfDefTema{С1.10}{S1-10}{Площадь поверхности и объём составного многогранника}{С1}{\tref{С1.1}}{3}{— нет в методичке}{}
\pfDefTema{С2.1}{S2-1}{Прямые в пространстве}{С2}{}{15}{стр. 120 (старое задание 14)}{}
\pfDefTema{С2.2}{S2-2}{Прямая и плоскость}{С2}{\tref{С2.1}}{15}{стр. 120–121 (старое задание 14)}{}
\pfDefTema{С2.3}{S2-3}{Две плоскости; параллельность плоскостей}{С2}{\tref{С2.2}}{15}{стр. 121–123 (старое задание 14)}{}
\pfDefTema{С2.4}{S2-4}{Перпендикулярность прямой и плоскости}{С2}{\tref{С2.2}}{15}{стр. 131–132 (старое задание 14)}{}
\pfDefTema{С2.5}{S2-5}{Теорема о трёх перпендикулярах}{С2}{\tref{С2.4}}{15}{стр. 132–133 (старое задание 14)}{}
\pfDefTema{С2.6}{S2-6}{Перпендикулярные плоскости}{С2}{\tref{С2.4}}{15}{стр. 124–125 (старое задание 14)}{}
\pfDefTema{С3.1}{S3-1}{Угол между прямыми}{С3}{\tref{С2.1}}{15}{стр. 123 (старое задание 14)}{}
\pfDefTema{С3.2}{S3-2}{Угол между прямой и плоскостью}{С3}{\tref{С2.4}}{15}{стр. 124 (старое задание 14)}{}
\pfDefTema{С3.3}{S3-3}{Угол между плоскостями}{С3}{\tref{С2.6}}{15}{стр. 124 (старое задание 14)}{}
\pfDefTema{С3.4}{S3-4}{Площадь ортогональной проекции}{С3}{\tref{С3.3}}{15}{стр. 125 (старое задание 14)}{}
\pfDefTema{С3.5}{S3-5}{Расстояние от точки до прямой}{С3}{\tref{П8.1}, \tref{П3.4}}{15}{стр. 126 (старое задание 14)}{}
\pfDefTema{С3.6}{S3-6}{Расстояние от точки до плоскости}{С3}{\tref{С2.4}, \tref{С2.5}}{15}{стр. 127–129 (старое задание 14)}{}
\pfDefTema{С3.7}{S3-7}{Расстояние между скрещивающимися прямыми}{С3}{\tref{С3.6}}{15}{стр. 129–130 (старое задание 14)}{}
\pfDefTema{С3.8}{S3-8}{Метод объёмов для расстояния от точки до плоскости}{С3}{\tref{С3.6}, \tref{С1.3}}{15}{— нет в методичке (обсуждается)}{}
\pfDefTema{С4.1}{S4-1}{Построение сечений, метод следов}{С4}{\tref{С2.3}}{15}{стр. 133 (старое задание 14)}{}
\pfDefTema{В1.1}{V1-1}{Перестановки, размещения, сочетания}{В1}{}{20}{стр. 213–215 (старое задание 19)}{}
\pfDefTema{В1.2}{V1-2}{Правило суммы и правило произведения}{В1}{}{4, 5, 20}{— нет в методичке}{}
\pfDefTema{В2.1}{V2-1}{Классическое определение вероятности}{В2}{\tref{Ч1.1}}{4, 5}{стр. 40–41 (старое задание 4–5)}{}
\pfDefTema{В2.2}{V2-2}{Свойства вероятности, противоположное событие}{В2}{\tref{В2.1}}{4, 5}{стр. 41 (старое задание 4–5)}{}
\pfDefTema{В2.3}{V2-3}{Сумма несовместных событий}{В2}{\tref{В2.2}}{4, 5}{стр. 44 (старое задание 4–5)}{}
\pfDefTema{В2.4}{V2-4}{Сумма совместных событий}{В2}{\tref{В2.3}}{4, 5}{стр. 44 (старое задание 4–5)}{}
\pfDefTema{В2.5}{V2-5}{Условная вероятность}{В2}{\tref{В2.1}}{4, 5}{стр. 44 (старое задание 4–5)}{}
\pfDefTema{В2.6}{V2-6}{Независимые события, произведение}{В2}{\tref{В2.5}}{4, 5}{стр. 43 (старое задание 4–5)}{}
\pfDefTema{В2.7}{V2-7}{Произведение зависимых событий}{В2}{\tref{В2.5}}{4, 5}{стр. 43 (старое задание 4–5)}{}
\pfDefTema{В2.8}{V2-8}{Формула полной вероятности, дерево вероятностей}{В2}{\tref{В2.7}, \tref{В2.3}}{5}{— нет в методичке}{}
\pfDefTema{В2.9}{V2-9}{Формула Бернулли}{В2}{\tref{В2.6}, \tref{В1.1}}{4, 5}{стр. 45 (старое задание 4–5)}{}
\pfDefTema{В3.1}{V3-1}{Случайная величина и её распределение}{В3}{\tref{В2.2}}{6}{— нет в методичке}{}
\pfDefTema{В3.2}{V3-2}{Математическое ожидание}{В3}{\tref{В3.1}}{6}{— нет в методичке}{}
\pfDefTema{В3.3}{V3-3}{Дисперсия и стандартное отклонение}{В3}{\tref{В3.2}}{6}{— нет в методичке}{}
\pfDefTema{В3.4}{V3-4}{Равномерное, показательное и нормальное распределения}{В3}{\tref{В3.3}}{6}{— нет в методичке (обсуждается)}{}
\pfDefTema{М1.1}{M1-1}{Работа с готовой формулой в прикладных задачах}{М1}{\tref{А5.2}}{10}{стр. 80 (старое задание 9)}{}
\pfDefTema{М1.2}{M1-2}{Движение по прямой}{М1}{\tref{Ч1.5}}{11}{стр. 86–87 (старое задание 10)}{}
\pfDefTema{М1.3}{M1-3}{Движение по воде}{М1}{\tref{М1.2}}{11}{стр. 87–88 (старое задание 10)}{}
\pfDefTema{М1.4}{M1-4}{Движение навстречу и вдогонку}{М1}{\tref{М1.2}}{11}{стр. 88–89 (старое задание 10)}{}
\pfDefTema{М1.5}{M1-5}{Работа и производительность}{М1}{}{11}{стр. 89–90 (старое задание 10)}{}
\pfDefTema{М1.6}{M1-6}{Смеси, сплавы, концентрация}{М1}{\tref{Ч1.3}}{11}{стр. 90–91 (старое задание 10)}{\rref{Р11-01}}
\pfDefTema{М1.7}{M1-7}{Средняя скорость}{М1}{\tref{М1.2}}{11}{— нет в методичке}{}
\pfDefTema{М2.1}{M2-1}{Многократное изменение на p\% (сложные проценты)}{М2}{\tref{Ч1.4}}{13}{стр. 152 (старое задание 16)}{}
\pfDefTema{М2.2}{M2-2}{Кредит: аннуитетные платежи}{М2}{\tref{М2.1}}{13}{стр. 152–153 (старое задание 16)}{}
\pfDefTema{М2.3}{M2-3}{Кредит: дифференцированные платежи}{М2}{\tref{М2.1}, \tref{А4.1}}{13}{стр. 153–155 (старое задание 16)}{}
\pfDefTema{М2.4}{M2-4}{Сравнение аннуитетной и дифференцированной схем}{М2}{\tref{М2.2}, \tref{М2.3}}{13}{стр. 155 (старое задание 16)}{}
\pfDefTema{М2.5}{M2-5}{Вклады}{М2}{\tref{М2.1}}{13}{стр. 155–156 (старое задание 16)}{}
\pfDefTema{М2.6}{M2-6}{Задачи на оптимизацию}{М2}{\tref{Д2.5}}{13, 17}{стр. 156–161 (старое задание 16)}{}

\pfDefRazbor{Р11-01}{R11-01}{11}{Смеси и сплавы}{\tref{М1.6}, \tref{Ч1.3}}{}
\pfDefRazbor{Р16-01}{R16-01}{16}{Неравенство с модулем}{\tref{А1.2}, \tref{А5.7}}{\tref{А5.8}}
\pfZadanieRazbory{11}{Р11-01}
\pfZadanieRazbory{16}{Р16-01}
